\section{Introduction}
\label{sec:introduction}

Scaling inference-time computation has emerged as a practical strategy for improving
the accuracy of large language models on reasoning tasks without additional model
training~\citep{snell2024scaling}.
By generating multiple candidate answers under a fixed inference budget and selecting
among them, a system can achieve accuracy gains comparable to those from model-parameter
scaling, at inference time~\citep{wang2023selfconsistency}.
This reframes answer quality as an \emph{intelligent decision problem}: given a pool of
already-generated candidates, which selection rule best exploits the available runtime
information under cost constraints?

In settings where the inference budget is \emph{matched per method}---every
candidate-producing method receives the same number of model calls per example---the
quality of the final answer depends entirely on the selection rule applied to the
already-generated pool.
Simple majority vote treats all candidates equally; yet a frontier generation run
produces structured metadata---records of search depth, branch structure, and override
events---that is richer than the raw candidates alone and may carry reliable signal
about candidate trustworthiness.
Exploiting such runtime-visible signals for dynamic selection is the core question
studied here.
This framing matches a practical deployment setting: after a service has generated a
fixed candidate pool under cost and rate-limit constraints, an auditable selector must
decide whether to trust the default frontier output or defer to external consensus,
without issuing extra inference calls or accessing ground-truth labels.

We examine this question in a concrete, reproducible setting.
The task is grade-school arithmetic reasoning on the GSM8K
benchmark~\citep{cobbe2021training}, the provider is Cohere command-r-plus-08-2024,
and the inference budget is $B = 6$ logical calls per example.
A frontier generator produces candidate answers via program-aided
execution~\citep{gao2023pal}; three independent external methods (L1, S1, and TALE)
provide reference answers under the same budget cap.
Given these already-generated outputs, the goal is to select a final answer without
additional inference calls and without using gold labels.

\paragraph{Contributions.}
The main contribution is \textbf{Failure-Trace Allocator (FTA)}, a deterministic
two-gate post-generation policy that uses failure-trace metadata logged during normal
frontier execution:
\begin{itemize}[leftmargin=1.4em,itemsep=2pt]
  \item The \textit{Low-Depth Guard} fires when the frontier's execution trace signals
        a structurally weak single-branch output, substituting an external majority-vote
        answer.
  \item The \textit{External-Consensus Fallback} fires when all three external baselines
        unanimously disagree with the frontier's selected answer.
\end{itemize}
Both gates are strictly gold-free: no correctness labels, ground-truth answers, or
example identifiers are accessed at decision time.
The two gates operate in strict priority order; at most one fires per example.

On the primary evaluation set of 300 disjoint examples (Final-300), FTA achieves
\textbf{86.67\%} accuracy.
On a disjoint Aggregate-720 combining three independent evaluation seeds (seeds 41,
61, and 71; $300 + 120 + 300$ examples with zero overlap in example identifiers or
question hashes), FTA achieves \textbf{80.69\%} accuracy versus 77.64\% (L1),
77.08\% (S1), and 75.14\% (TALE).
Source-stratified 95\% bootstrap confidence intervals are strictly positive for all
four pairwise comparisons, including versus the best external source
(CI lower bound $+1.11$~pp).
As supporting evidence, a multi-provider regime-learning selector (D9) trained on the
same failure-trace signal improves over the frontier baseline in cross-validation with zero
false-override events.
FTA remains ahead of a pooled four-answer ensemble (frontier+L1+S1+TALE) by point
estimate on both evaluation sets, but that comparison is not statistically separated.

We also document five not-retained prototype variants evaluated against the same
criteria; these are reported as negative or insufficient results in
Section~\ref{sec:ablation} and Table~\ref{tab:negative_results}.

\paragraph{Scope disclosures.}
Four scope notes are integral to interpreting all claims in this paper.
\emph{(i)} Confidence intervals are source-stratified bootstrap; the interval versus
the pooled four-answer ensemble includes zero on both Final-300 and Aggregate-720,
so superiority over ensembling is \emph{not} claimed.
\emph{(ii)} The matched-budget contract is an evaluation design: constructing the full
frontier+L1+S1+TALE pool from scratch corresponds to 24 logical calls per example,
not six.
\emph{(iii)} All evaluation and selection findings are scoped to GSM8K with the
Cohere command-r-plus-08-2024 provider; generalization to other benchmarks or
providers requires independent replication.
\emph{(iv)} The 120-example evaluation set at seed 61 is failure-enriched by
construction; results on seeds 41 and 71 serve as out-of-distribution checks.

\paragraph{Paper organization.}
Section~\ref{sec:related_work} reviews related work on test-time compute scaling,
verifier-guided reranking, selective prediction, and cost-aware routing.
Section~\ref{sec:method} specifies FTA precisely, including exact trigger conditions
and a policy summary table.
Sections~\ref{sec:experiments} and~\ref{sec:results} describe the evaluation protocol
and report results with full confidence intervals.
Section~\ref{sec:ablation} summarizes ablation variants and not-retained outcomes.
Sections~\ref{sec:discussion}--\ref{sec:conclusion} discuss implications, limitations,
and directions for future work.
